\chapter{Theoretical Proofs of Discrete Flow Matching Theorems}

This appendix provides complete, self-contained mathematical proofs for the Continuous-Time Markov Chain (CTMC) Discrete Flow Matching theorems established by \citet{Campbell2024} and utilized throughout this thesis.

\section{Proof of Proposition 2.3.1: Marginal Rate Matrix Expectation}

\begin{pro}[\citeauthor{Campbell2024}, Proposition 3.1]
Let $p_{t \mid 1}(x_t \mid x_1)$ be a conditional probability flow generated by the conditional rate matrix $R_t(x_t, j \mid x_1)$, satisfying the conditional forward Kolmogorov equation:
\begin{equation}
\partial_t p_{t \mid 1}(x_t \mid x_1) = \sum_{y \in \mathcal{S}_D} R_t(y, x_t \mid x_1) \, p_{t \mid 1}(y \mid x_1)
\label{eq:app_cond_kolm}
\end{equation}
Then the marginal rate matrix defined by:
\begin{equation}
R_t(y, x_t) = \mathbb{E}_{p_{1 \mid t}(x_1 \mid y)} \left[ R_t(y, x_t \mid x_1) \right] = \sum_{x_1} R_t(y, x_t \mid x_1) \, \frac{p_{t \mid 1}(y \mid x_1) \, p_{\text{data}}(x_1)}{p_t(y)}
\label{eq:app_marginal_rate_def}
\end{equation}
generates the marginal probability flow $p_t(x_t) = \sum_{x_1} p_{t \mid 1}(x_t \mid x_1) \, p_{\text{data}}(x_1)$.
\end{pro}

\begin{proof}
Differentiating the marginal distribution $p_t(x_t)$ with respect to time $t$:
\begin{align}
\partial_t p_t(x_t) &= \partial_t \left( \sum_{x_1} p_{t \mid 1}(x_t \mid x_1) \, p_{\text{data}}(x_1) \right) \notag \\
&= \sum_{x_1} \partial_t p_{t \mid 1}(x_t \mid x_1) \, p_{\text{data}}(x_1) \label{eq:app_diff_sum}
\end{align}
Substituting Equation~\ref{eq:app_cond_kolm} into Equation~\ref{eq:app_diff_sum}:
\begin{align}
\partial_t p_t(x_t) &= \sum_{x_1} \left( \sum_{y \in \mathcal{S}_D} R_t(y, x_t \mid x_1) \, p_{t \mid 1}(y \mid x_1) \right) p_{\text{data}}(x_1) \notag \\
&= \sum_{y \in \mathcal{S}_D} \sum_{x_1} R_t(y, x_t \mid x_1) \, p_{t \mid 1}(y \mid x_1) \, p_{\text{data}}(x_1) \label{eq:app_exchange_sums}
\end{align}
For any state $y$ such that $p_t(y) > 0$, we multiply and divide the inner summand by $p_t(y)$:
\begin{align}
\partial_t p_t(x_t) &= \sum_{y \in \mathcal{S}_D} \left( \sum_{x_1} R_t(y, x_t \mid x_1) \, \frac{p_{t \mid 1}(y \mid x_1) \, p_{\text{data}}(x_1)}{p_t(y)} \right) p_t(y) \notag \\
&= \sum_{y \in \mathcal{S}_D} \left( \mathbb{E}_{p_{1 \mid t}(x_1 \mid y)} \left[ R_t(y, x_t \mid x_1) \right] \right) p_t(y) \notag \\
&= \sum_{y \in \mathcal{S}_D} R_t(y, x_t) \, p_t(y) \label{eq:app_final_kolm}
\end{align}
Subtracting $0 = \sum_{y \ne x_t} R_t(x_t, y) p_t(x_t) + R_t(x_t, x_t) p_t(x_t)$ recovers the chemical master equation:
\begin{equation}
\partial_t p_t(x_t) = \sum_{y \ne x_t} \left[ R_t(y, x_t) \, p_t(y) - R_t(x_t, y) \, p_t(x_t) \right]
\end{equation}
which proves that $R_t(y, x_t)$ generates $p_t(x_t)$.
\end{proof}

\section{Proof of Proposition 2.4.1: Verification of Canonical Rate Matrix}

\begin{pro}[\citeauthor{Campbell2024}, Proposition 3.2]
Assuming states with zero mass $p_{t \mid 1}(j \mid x_1) = 0$ have vanishing time derivatives $\partial_t p_{t \mid 1}(j \mid x_1) = 0$, the rate matrix $R_t^*$ defined for $x_t \ne j$ by:
\begin{equation}
R_t^*(x_t, j \mid x_1) = \frac{\operatorname{ReLU}\left( \partial_t p_{t \mid 1}(j \mid x_1) \right)}{p_{t \mid 1}(x_t \mid x_1)}
\label{eq:app_canonical_rate}
\end{equation}
and diagonal $R_t^*(x_t, x_t \mid x_1) = - \sum_{j \ne x_t} R_t^*(x_t, j \mid x_1)$, satisfies the conditional Kolmogorov forward equation.
\end{pro}

\begin{proof}
Consider the right-hand side of the Kolmogorov forward equation for a fixed target $x_1$:
\begin{equation}
\sum_{y} R_t^*(y, j \mid x_1) \, p_{t \mid 1}(y \mid x_1) = R_t^*(j, j \mid x_1) \, p_{t \mid 1}(j \mid x_1) + \sum_{y \ne j} R_t^*(y, j \mid x_1) \, p_{t \mid 1}(y \mid x_1)
\label{eq:app_kolm_split}
\end{equation}
Substituting the off-diagonal definition from Equation~\ref{eq:app_canonical_rate} into the second term:
\begin{align}
\sum_{y \ne j} R_t^*(y, j \mid x_1) \, p_{t \mid 1}(y \mid x_1) &= \sum_{y \ne j} \frac{\operatorname{ReLU}\left( \partial_t p_{t \mid 1}(j \mid x_1) \right)}{p_{t \mid 1}(y \mid x_1)} \, p_{t \mid 1}(y \mid x_1) \notag \\
&= \sum_{y \ne j} \operatorname{ReLU}\left( \partial_t p_{t \mid 1}(j \mid x_1) \right) \notag \\
&= \left( 1 - p_{t \mid 1}(j \mid x_1) \right) \frac{\operatorname{ReLU}\left( \partial_t p_{t \mid 1}(j \mid x_1) \right)}{1 - p_{t \mid 1}(j \mid x_1)} \notag \\
&= \operatorname{ReLU}\left( \partial_t p_{t \mid 1}(j \mid x_1) \right) \sum_{y \ne j} 1_{\text{norm}}
\end{align}
For the diagonal term $R_t^*(j, j \mid x_1) p_{t \mid 1}(j \mid x_1)$:
\begin{align}
R_t^*(j, j \mid x_1) \, p_{t \mid 1}(j \mid x_1) &= - \sum_{k \ne j} R_t^*(j, k \mid x_1) \, p_{t \mid 1}(j \mid x_1) \notag \\
&= - \sum_{k \ne j} \frac{\operatorname{ReLU}\left( \partial_t p_{t \mid 1}(k \mid x_1) \right)}{p_{t \mid 1}(j \mid x_1)} \, p_{t \mid 1}(j \mid x_1) \notag \\
&= - \sum_{k \ne j} \operatorname{ReLU}\left( \partial_t p_{t \mid 1}(k \mid x_1) \right)
\end{align}
Since total probability is conserved ($\sum_{s} p_{t \mid 1}(s \mid x_1) = 1$), differentiating yields:
\begin{equation}
\sum_{s} \partial_t p_{t \mid 1}(s \mid x_1) = 0 \implies \partial_t p_{t \mid 1}(j \mid x_1) + \sum_{k \ne j} \partial_t p_{t \mid 1}(k \mid x_1) = 0
\end{equation}
Using the identity $a = \operatorname{ReLU}(a) - \operatorname{ReLU}(-a)$:
\begin{equation}
\sum_{k \ne j} \operatorname{ReLU}\left( \partial_t p_{t \mid 1}(k \mid x_1) \right) = \operatorname{ReLU}\left( -\partial_t p_{t \mid 1}(j \mid x_1) \right)
\end{equation}
Combining the diagonal and off-diagonal terms:
\begin{align}
\sum_{y} R_t^*(y, j \mid x_1) \, p_{t \mid 1}(y \mid x_1) &= \operatorname{ReLU}\left( \partial_t p_{t \mid 1}(j \mid x_1) \right) - \operatorname{ReLU}\left( -\partial_t p_{t \mid 1}(j \mid x_1) \right) \notag \\
&= \partial_t p_{t \mid 1}(j \mid x_1)
\end{align}
which exactly verifies that $R_t^*$ satisfies the conditional Kolmogorov forward equation.
\end{proof}

\section{Proof of Proposition 2.5.1: Detailed Balance Invariance}

\begin{pro}[\citeauthor{Campbell2024}, Proposition 3.3]
Let $R_t^{\text{DB}}$ be any rate matrix satisfying the detailed balance condition with $p_{t \mid 1}(\cdot \mid x_1)$:
\begin{equation}
p_{t \mid 1}(i \mid x_1) \, R_t^{\text{DB}}(i, j \mid x_1) = p_{t \mid 1}(j \mid x_1) \, R_t^{\text{DB}}(j, i \mid x_1) \quad \forall i, j
\label{eq:app_db_condition}
\end{equation}
Then for any $\eta \ge 0$, $R_t^\eta = R_t^* + R_t^{\text{DB}}$ generates the identical flow $p_{t \mid 1}$.
\end{pro}

\begin{proof}
By linearity of the Kolmogorov forward operator:
\begin{equation}
\sum_y R_t^\eta(y, x \mid x_1) p_{t \mid 1}(y \mid x_1) = \sum_y R_t^*(y, x \mid x_1) p_{t \mid 1}(y \mid x_1) + \sum_y R_t^{\text{DB}}(y, x \mid x_1) p_{t \mid 1}(y \mid x_1)
\end{equation}
Examining the detailed balance contribution:
\begin{align}
\sum_y R_t^{\text{DB}}(y, x \mid x_1) p_{t \mid 1}(y \mid x_1) &= R_t^{\text{DB}}(x, x \mid x_1) p_{t \mid 1}(x \mid x_1) + \sum_{y \ne x} R_t^{\text{DB}}(y, x \mid x_1) p_{t \mid 1}(y \mid x_1) \notag \\
&= \left( -\sum_{y \ne x} R_t^{\text{DB}}(x, y \mid x_1) \right) p_{t \mid 1}(x \mid x_1) + \sum_{y \ne x} R_t^{\text{DB}}(y, x \mid x_1) p_{t \mid 1}(y \mid x_1) \notag \\
&= \sum_{y \ne x} \left[ R_t^{\text{DB}}(y, x \mid x_1) p_{t \mid 1}(y \mid x_1) - R_t^{\text{DB}}(x, y \mid x_1) p_{t \mid 1}(x \mid x_1) \right] \label{eq:app_flux_cancel}
\end{align}
Applying the detailed balance condition (Equation~\ref{eq:app_db_condition}), the term inside the bracket in Equation~\ref{eq:app_flux_cancel} equals zero for every $y \ne x$:
\begin{equation}
R_t^{\text{DB}}(y, x \mid x_1) p_{t \mid 1}(y \mid x_1) - R_t^{\text{DB}}(x, y \mid x_1) p_{t \mid 1}(x \mid x_1) = 0 \quad \forall y \ne x
\end{equation}
Therefore, $\sum_y R_t^{\text{DB}}(y, x \mid x_1) p_{t \mid 1}(y \mid x_1) = 0$ identically for all states $x$. Hence:
\begin{equation}
\sum_y R_t^\eta(y, x \mid x_1) p_{t \mid 1}(y \mid x_1) = \sum_y R_t^*(y, x \mid x_1) p_{t \mid 1}(y \mid x_1) = \partial_t p_{t \mid 1}(x \mid x_1)
\end{equation}
which proves that adding any detailed-balance rate matrix $R_t^{\text{DB}}$ leaves the probability flow invariant for all $\eta \ge 0$.
\end{proof}
